\documentclass[12pt]{article}
\usepackage{auck2e,multicol,amssymb,amsmath,color}
\usepackage{epsfig}
\usepackage[english]{babel}
\usepackage{blindtext}
\usepackage{microtype}
\usepackage{minted}

%\usepackage{graphicx}
\newcommand{\ds}{\displaystyle}
\newcommand{\mmark}[1]{[{#1 mark(s)}]}
\newcommand{\markk}[1]{[{\bf #1 marks}]}
\newcommand{\dspace}{\renewcommand{\baselinestretch}{2} \small \normalsize}
\newcommand{\sspace}{\renewcommand{\baselinestretch}{1} \small \normalsize}
\newcommand{\btheta}{\mbox{\boldmath $\theta$}}
\newcommand{\blambda}{\mbox{\boldmath $\lambda$}}
\newcommand{\bTheta}{\mbox{\boldmath $\Theta$}}
\newcommand{\bzeta}{\mbox{\boldmath $\zeta$}}
\newcommand{\bpi}{\mbox{\boldmath $\pi$}}
\newcommand{\bphi}{\mbox{\boldmath $\phi$}}
\newcommand{\ba}{\mbox{\boldmath $a$}}
\newcommand{\bb}{\mbox{\boldmath $b$}}
\newcommand{\bd}{\mbox{\boldmath $d$}}
\newcommand{\bm}{\mbox{\boldmath $m$}}
\newcommand{\bx}{\mbox{\boldmath $x$}}
\newcommand{\bh}{\mbox{\boldmath $h$}}
\newcommand{\bu}{\mbox{\boldmath $u$}}
\newcommand{\bP}{\mbox{\boldmath $P$}}
\newcommand{\bQ}{\mbox{\boldmath $Q$}}
\newcommand{\bR}{\mbox{\boldmath $R$}}
\newcommand{\by}{\mbox{\boldmath $y$}}
\newcommand{\bY}{\mbox{\boldmath $Y$}}
\newcommand{\bX}{\mbox{\boldmath $X$}}
\newcommand{\bU}{\mbox{\boldmath $U$}}
\newcommand{\bD}{\mbox{\boldmath $D$}}
\newcommand{\bV}{\mbox{\boldmath $V$}}
\newcommand{\bH}{\mbox{\boldmath $H$}}
\newcommand{\bgamma}{\mbox{\boldmath $\gamma$}}
\newcommand{\bGamma}{\mbox{\boldmath $\Gamma$}}
\newcommand{\balpha}{\mbox{\boldmath $\alpha$}}
\newcommand{\bbeta}{\mbox{\boldmath $\beta$}}
\newcommand{\bdelta}{\mbox{\boldmath $\delta$}}
\newcommand{\bomega}{\mbox{\boldmath $\omega$}}
\newcommand{\bpsi}{\mbox{\boldmath $\psi$}}
\newcommand{\bepsilon}{\mbox{\boldmath $\epsilon$}}
\newcommand{\be}{\begin{eqnarray}}
\newcommand{\ee}{\end{eqnarray}}
\def\BackgroundEPS#1#2#3#4{%
    \special{ps: @beginspecial @setspecial initmatrix
                0.9 setgray #2 #3 translate #4 dup scale}
    \special{ps: plotfile #1}
    \special{ps: @endspecial}
}
\newcommand{\wbugs}{{\sffamily\textbf{WinBUGS }}}
\newcommand{\coda}{{\sffamily\textbf{CODA }}}
\newcommand{\glm}{{\sffamily\textbf{G.L.M }}}
\newcommand{\rwb}{{\sffamily\textbf{R2WinBUGS }}}
\newcommand{\rar}{{\sffamily\textbf{Rock and roll man! }}}
\newcommand{\anova}{{\sffamily\textbf{ANOVA }}}
\newcommand{\mcmc}{{\sffamily\textbf{MCMC }}}
\newcommand{\dic}{{\sffamily\textbf{DIC }}}
\newcommand{\dagr}{{\sffamily\textbf{DAG }}}
\newcommand{\boxeq}[1]{%
\fboxrule=0.5mm%
\fboxsep=5mm%
\begin{center}\fcolorbox{black}{yellow}{\begin{minipage}{0.9\linewidth}
\begin{Beqnarray} #1 \end{Beqnarray}\end{minipage}}\end{center}%
\fboxrule=1pt%
\fboxsep=3pt%
}

\newcommand{\boxst}[1]{%
\fboxrule=0.5mm%
\fboxsep=5mm%
\begin{center}\fcolorbox{black}{yellow}{ \begin{minipage}{0.9\linewidth}#1 \end{minipage} }\end{center}%
\fboxrule=1pt%
\fboxsep=3pt%
}
%%%%%%%%%%%%%%%%%%%%%%%%%%%%%%%%%%%%%%%%%%%%%%%%%%%%%%%%%%%%%%%%%%%%%%%%
% Macro         :   \BackgroundEPS{filename}{tx}{ty}{s}
%-----------------------------------------------------------------------
% Description   :   Display an EPS file in the background of the
%                   current page.
%-----------------------------------------------------------------------
% Author        :   Bernard Desruisseaux (bernard@XYZZYcs.mcgill.ca)
%-----------------------------------------------------------------------
% Usage         :   \BackgroundEPS{pic.eps}{80}{100}{2}
%                   \BackgroundEPS{"`gunzip -c pic.eps.gz"}{80}{100}{2}
%                   \BackgroundEPS{"`zcat pic.eps.Z"}{80}{100}{2}
%%%%%%%%%%%%%%%%%%%%%%%%%%%%%%%%%%%%%%%%%%%%%%%%%%%%%%%%%%%%%%%%%%%%%%%%

\newcommand{\given}{\,|\,}

\paper{Introduction to Bayesian Statistics}
\due{Not for assessment}
\thetitle{\textbf{Tutorial 3}} 
\begin{document}


%\BackgroundEPS{bayes.eps}{0}{0}{1.2}
\makeband


\begin{Questions}
\item The file {\tt american\_football.R} contains measurements of the weights
of players on five NFL teams, measured in kilograms. 17 players were measured
from each team. You will apply the Bayesian ``one-way ANOVA'' model from lectures
to this data.

\begin{enumerate}
\item Make boxplots of the data using R's {\tt boxplot} function.
\item Write up the Bayesian ``one-way ANOVA'' model and modify the
first {\tt for} loop so there are five {\tt mu} parameters instead of four.
\item According to the assumptions in the model, what is the prior probability
of the hypothesis $H_0: \mu_1 = \mu_2 = \mu_3 = \mu_4 = \mu_5$? What will the
posterior probability of $H_0$ be?
\item Run the code and inspect the trace plots.
\item The data suggest Team 3 (Denver Broncos)
is the heaviest and Team 5 (San Francisco 49ers) is the lightest, on average.
Find the posterior probability that $\mu_3 > \mu_5$.
\item Look at the posterior distribution for the {\tt log\_diversity}
hyperparameter. What does it show?
\end{enumerate}


\item
The file {\tt kobe\_bryant.csv} contains data on the number of shots
(field goals) attempted by basketball player Kobe Bryant, and the number of successes achieved, for each season of his NBA career.
The only two columns you will use are attempts and successes. Use the following hierarchical
model:
\begin{minted}{r}
model
{
    grand_mean ~ dunif(-10, 10)
    diversity ~ dunif(0, 3)
    for(i in 1:length(successes))
    {
        l[i] ~ dnorm(grand_mean, 1/diversity^2)
        logit(theta[i]) <- l[i]
        successes[i] ~ dbin(theta[i], attempts[i])
    }
}
\end{minted}

The logit transformation is used because a normal prior would allow values
outside the range $[0, 1]$. The logit transformation makes sure that
the success probability for each season, $\theta$, is between 0 and 1.

\begin{enumerate}
\item[(a)] Identify the hyperparameters, the parameters, and the data.
\item[(b)] How greatly did Kobe Bryant’s shooting ability vary across seasons?
\item[(c)] Can you think of any interesting changes you could make to the model? If you
can, try implementing them.
\item[(d)] Instead of the hierarchical model, try analysing each season's
\mintinline{r}{attempts} and \mintinline{r}{successes} using a Uniform$(0, 1)$
prior for the success probability $\theta$ for each season. Analytically
compute the posterior mean and standard deviation using the formulas
for the Beta distribution (you can find these on Wikipedia or elsewhere).
\item[(e)] Compare the posterior means and standard deviations from (d)
to what you get from the hierarchical model.
\end{Questions}

\end{document}




